% INITIAL WRITE-UP DRAFT WITH ORIGINAL TEXT AND PLACEHOLDERS
% Preserved in project_history/ for reference and technical provenance.

\chapter{Introduction}

\section{Background of the Study}
Robotics has moved beyond isolated industrial automation toward shared, human-centered environments. In these settings, robots are expected to operate intelligently alongside people rather than executing rigid, isolated commands. This shift has made Human-Robot Interaction (HRI) a major area of research. For robots to function effectively, they must interpret human behavior in a way that is safe, natural, and responsive. 

One critical aspect of this is human intention recognition through hand gestures. While recent developments in computer vision have made vision-based gesture recognition more practical, many existing systems rely heavily on cloud computing or high-power graphics processors (GPUs) \cite{mahmud2022}. Furthermore, converting visual data into reliable robotic action in real-time remains challenging, especially on low-power edge computing devices where reaction time thresholds are critical \cite{tsitos2022}. A robotic system must not only recognize a gesture but also understand its own spatial environment to execute the command safely. This study addresses these challenges by developing a vision-based robotic application that integrates gesture recognition and autonomous navigation directly on an edge computing platform.

\section{Problem Statement}
Despite significant advancements in collaborative robotics, achieving natural and reliable HRI remains a challenge, particularly for systems operating on resource-constrained hardware. Most commercial and research robots still rely on rigid control paradigms requiring physical interfaces, touchscreens, or voice commands, which are often impractical in noisy industrial settings or sterile healthcare environments. 

Furthermore, existing vision-based gesture recognition systems typically present three major engineering bottlenecks. First, they often rely heavily on cloud infrastructure or high-power GPU acceleration, making them unsuitable for edge deployment where connectivity is unreliable \cite{mahmud2022}. Second, there is frequently a feature mismatch between the spatial data used during model training and the processed data available during live deployment, leading to classification failures in production. Finally, many gesture-controlled robots act blindly; they process commands without spatial awareness of their environment, creating safety hazards. 

There is a critical need for a lightweight, edge-deployable system that solves these issues. This study addresses the lack of a unified, GPU-independent framework that integrates real-time hand gesture recognition and autonomous navigation within a robust robotic middleware architecture.

\section{Objectives of the Study}
\subsection{General Objective}
The general objective of this study is to design, implement, and evaluate an edge-deployed, real-time human intention recognition system that integrates hand gesture classification and autonomous navigation into a unified ROS2 middleware pipeline on a mobile robotic platform.

\subsection{Specific Objectives}
\begin{enumerate}
    \item To design and implement a real-time hand gesture recognition pipeline utilizing vision-based landmark tracking and a Multilayer Perceptron (MLP) classifier based on engineered geometric hand features.
    \item To develop and train a static gesture classification model achieving measurable accuracy and reliability using a custom-collected, balanced dataset.
    \item To implement a ROS2-based multi-node software architecture integrating gesture recognition and centralised decision-making for autonomous robot control.
    \item To integrate Simultaneous Localization and Mapping (SLAM) and autonomous navigation algorithms to enable gesture-directed robot operation within a mapped environment.
    \item To evaluate the performance of the deployed system using classification metrics, system latency, and real-time hardware throughput measurements on an edge computing platform.
    \item To develop and integrate a sequential motion prediction model capable of anticipating an operator's short-horizon movement from pose landmark data, complementing explicit gesture commands with implicit motion-based intention inference.
\end{enumerate}

\section{Scope of the Study}
The scope of this research is bounded to the development of a real-time, gesture-controlled navigation system for an autonomous mobile robot. The operational environment is restricted to structured, indoor settings where lighting and occupancy are controlled to ensure predictable sensor performance. To ensure safe and predictable interaction, the human-robot interface is limited strictly to the visual recognition of predefined, static hand gestures, and the system is designed to respond to only a single active operator at any given time. Furthermore, to address the constraints of network dependency and latency, the project focuses on localized processing; the entire perception and navigation architecture is capable of operating using only the onboard computational resources of the robot, without reliance on remote cloud infrastructure. "Real-time" operation is defined as the system’s ability to maintain a processing cycle capable of supporting responsive interaction at human-speed.

\section{Significance of the Study}
The significance of this study lies in its approach to decoupling human intention recognition from cloud infrastructure. By bridging vision-based gesture classification with spatial navigation directly on the robot, this research provides a scalable framework for autonomous control in resource-constrained environments. 

Practically, a system capable of interpreting natural hand gestures without physical contact addresses specific industrial and healthcare needs. In acoustically complex environments, it provides a reliable control alternative where voice commands fail. In healthcare settings, it facilitates touch-less interaction, maintaining strict sterility protocols. Academically, this study contributes a modular, edge-deployed framework for multi-modal perception, demonstrating how machine learning inference can be optimized for low-power hardware without sacrificing classification reliability. Ultimately, this work seeks to promote safer and more accessible human-robot collaboration through localized, low-power computational frameworks, ensuring operational predictability in accordance with international safety standards for collaborative robotic workspaces \cite{iso150662016}.

\section{Organization of the Study}
This thesis is structured into five chapters. Chapter One introduces the research background, problem statement, objectives, scope, and significance. Chapter Two presents a review of related literature concerning human intention recognition, edge-based machine learning, and robotic navigation. Chapter Three details the system methodology, including hardware specifications, software architecture, model training procedures, and middleware integration. Chapter Four presents the testing protocols, performance metrics, and an analysis of the system's operational capabilities. Finally, Chapter Five concludes the study by summarizing the findings, outlining system constraints, and suggesting recommendations for future work.

% ================= CHAPTER TWO =================
\cleardoublepage

\chapter{Literature Review}

\section*{2.0 Overview}
\addcontentsline{toc}{section}{2.0 Overview}

This chapter reviews previous studies on human intention recognition, gesture and motion analysis, and robot navigation in human--robot interaction. Existing research shows that robots can recognize gestures, predict human actions from motion, and improve navigation by considering obstacle behavior, contributing to safer and more efficient collaboration. However, many of these systems operate under specific conditions, rely on limited sensing approaches, or treat motion recognition, gesture recognition, and navigation as separate components, which can affect reliability in practical environments. Therefore, this study builds on these existing approaches by seeking to improve intention recognition through the integration of human motion and hand gesture analysis, with the aim of developing a system capable of recognizing and predicting human intentions more reliably, even under varying conditions in a mapped environment.

% ================= THEORETICAL REVIEW =================
\section{Theoretical Review}

This study is theoretically grounded in concepts from Human--Robot Interaction (HRI), human intention theory, motion and gesture recognition, and spatial navigation. Human intention theory posits that human actions are goal-directed and that early segments of motion contain predictive information about future behavior, forming the foundation for anticipatory robotic systems. Within Human--Robot Interaction theory, modern collaborative robots are expected to move beyond reactive responses toward proactive and context-aware behavior in shared environments.

The theoretical basis of gesture and motion recognition in this study relies on the
assumption that human hand configurations and short-term movement trajectories follow
patterns that can be modeled geometrically and temporally. Gesture classification is
grounded in the assumption that a hand's spatial landmark configuration at a single
instant carries sufficient information to distinguish a fixed vocabulary of commands,
supporting a feed-forward classification approach. Movement prediction, by contrast,
requires temporal context — a single frame cannot reveal direction or intent —
motivating the use of a sequential model (Long Short-Term Memory) capable of learning
short-horizon motion patterns from consecutive pose observations.

Spatial mapping and localization theory provides the framework for autonomous
navigation between human encounters, using occupancy-grid-based Simultaneous
Localization and Mapping (SLAM) and probabilistic localization to allow the robot to
safely traverse a mapped environment independently of human interaction events.

Finally, multi-modal fusion theory justifies the integration of human motion and hand gesture cues, asserting that combining multiple sensory modalities enhances robustness, reduces ambiguity, and improves predictive reliability. Together, these theoretical foundations support the development of an integrated system capable of recognizing and predicting human intentions within a mapped, physically deployed robotic environment.

\section{Review of Related Works}

\subsection{Review One}

\subsubsection{Research Topic and Author}

Real-time Feasibility of a Human Intention Method Evaluated Through a Competitive Human--Robot Reaching Game (Tsitos et al., 2022)

\subsubsection{Objectives}

The main objective of the study was to develop and evaluate a real-time human intention prediction system using a single RGB-D camera and OpenPose, integrate this prediction module into a ROS-based robot control pipeline, and determine whether the robot could act quickly enough based on early intention prediction. The system was evaluated through a competitive human--robot reaching game.

\subsubsection{Methodology}

The study adopted a ROS-integrated machine learning approach combined with real-time human pose estimation. An RGB-D camera was used to capture human motion, while OpenPose was applied to extract wrist joint positions. The data were pre-processed to remove noise and detect movement phases. Feature vectors were constructed and used to train classifiers including Logistic Regression, Support Vector Machine, Decision Tree, and Naive Bayes. The trained model was integrated into a ROS pipeline to control a UR3 robot in real time.

\subsubsection{Strengths}

The study demonstrates that accurate intention prediction can be achieved using a simple and cost-effective setup with a single RGB-D camera. It supports real-time prediction using minimal data points and integrates seamlessly into a ROS framework. The use of a real robotic system enhances practical relevance, and the identification of a reaction time threshold highlights the importance of system latency.

\subsubsection{Limitations}

The study did not evaluate scenarios involving direct interaction between the human and robot in the same workspace. As a result, collision avoidance and motion adaptation were not tested, limiting its applicability in real collaborative environments.

\subsubsection{Research Gap}

Tsitos et al. demonstrate that intention prediction can be operationally fast using a
minimal wrist-position feature set, but their evaluation is confined to a stationary
robotic arm reaching toward a fixed target — the robot and human never share a mobile
workspace, no autonomous navigation is involved, and the only "communication channel"
is inferred motion, with no explicit gesture command available to the human operator.
This leaves a gap between fast intention prediction and its integration into a mobile,
autonomously navigating platform capable of also accepting deliberate human commands —
precisely the combination this project implements.

\subsection{Review Two}

\subsubsection{Research Topic and Author}

3D Gesture Recognition and Adaptation for Human--Robot Interaction (J. A. Mahmud et al., 2022)

\subsubsection{Objectives}

The objective of the study was to develop a Kinect-based three-dimensional hand gesture recognition and adaptation system capable of recognizing pointing and dynamic gestures in real time and applying them to control robot navigation within human--robot interaction environments.

\subsubsection{Methodology}

The study employed a Kinect-based skeletal tracking system to capture three-dimensional joint coordinates of the user’s hands and body. Pointing direction was estimated using geometric line equations on the X--Z plane, while dynamic gestures were identified from sequences of joint movements over time. The captured coordinates were translated and normalized using a fixed reference point to ensure consistency in feature representation. These features were used to train and evaluate three machine learning classifiers, namely Hidden Markov Model, Support Vector Machine, and Convolutional Neural Network.

In addition, a gesture adaptation module was introduced to allow the system to learn new gestures through user consent and semi-supervised learning. Robot navigation based on recognized gestures was implemented within a simulation environment, and the system was trained and tested using a dataset of 3,600 gesture samples collected from 24 participants across different age groups.

\subsubsection{Strengths}

The study demonstrates several strengths in both design and implementation. The use of three-dimensional joint coordinate data enhances robustness against variations in lighting conditions and background complexity compared to two-dimensional approaches. The integration of both pointing and dynamic gestures enables richer and more natural human--robot interaction. The system supports real-time gesture recognition, making it suitable for interactive applications.

Furthermore, the inclusion of a gesture adaptation mechanism improves system flexibility by allowing the incorporation of new gestures over time. The comparative evaluation of multiple machine learning models provides insight into optimal model performance. The system was also tested across different age groups, demonstrating its adaptability to diverse users. In addition, the development of a graphical user interface and simulation platform facilitates user interaction and system evaluation. The study also proposes structured navigation strategies based on recognized gestures.

\subsubsection{Limitations}

Despite its strengths, the system presents several limitations. It relies entirely on the Kinect sensor, meaning its performance is dependent on the sensor’s field of view and skeletal tracking accuracy, and it is limited to tracking a single user at a time. Gesture recognition requires users to maintain specific postures, which may not reflect natural human behavior. Furthermore, robot navigation was evaluated only in a simulated environment, meaning real-world uncertainties were not considered.

The system was initially trained using a limited set of predefined gestures, and the process of learning new gestures requires repeated user interaction and system retraining. In addition, variations in performance across different age groups were observed, and the assumption of fixed spatial constraints may limit the system’s generalization to other environments.

\subsubsection{Research Gap}

The system estimates pointing direction using only right-hand joint data, specifically the hand, elbow, and shoulder, which restricts its ability to capture more natural and complex human pointing behaviors involving both hands or finer finger movements. Additionally, robot navigation is based solely on gesture direction without incorporating obstacle detection or avoidance mechanisms. This reduces the system’s safety and practicality in real-world environments. These limitations highlight the need for a more robust gesture interpretation approach integrated with safe navigation strategies that account for environmental obstacles and human presence.

While Mahmud et al.'s system and the present work both rely on visual sensing,
Zafar et al.~\cite{zafar2023} demonstrate an alternative, non-visual approach: a
hybrid deep learning model for hand gesture recognition using surface
electromyography (sEMG) sensors. This sensing modality sidesteps the lighting- and
occlusion-sensitivity inherent to camera-based methods entirely, at the cost of
requiring the operator to wear dedicated sensing hardware -- a trade-off directly
opposite to the device-free, camera-based interaction model this project and Mahmud
et al. both prioritize.

\subsection{Review Three}

\subsubsection{Research Topic and Author}

Safe and Efficient Motion Planning for Material Transportation Robots considering Intention Prediction of Obstacles (Y. Li, H. Zhang et al., 2023)

\subsubsection{Objectives}

The main objective of the study was to develop a motion planning approach for material transportation robots that improves both safety and efficiency in construction sites by extending beyond traditional obstacle avoidance methods. Specifically, the study aimed to enable robots to predict the movement intention of obstacles, such as determining whether a worker is likely to move away, and incorporate this prediction into the navigation process through costmap updates without modifying existing planning algorithms. The approach also sought to reduce unnecessary obstacle avoidance while maintaining safe and smooth navigation in dynamic human–robot shared environments.

\subsubsection{Methodology}

The study proposed an intention-aware navigation framework for material transportation robots operating in construction environments. The system combined camera-based object detection and LiDAR sensing to identify and classify different types of obstacles, including workers, carts, machines, and materials. Instead of treating all obstacles equally, the system predicted whether an obstacle was likely to actively move away from the robot based on its type and distance using a ReLU-based intention model.

This predicted intention was used to modify the robot’s global costmap for path planning, while the local costmap continued to update normally to ensure safety. The updated costmaps were integrated into the Nav2 ROS2 navigation framework without altering the original planning algorithm. The approach was validated in Isaac Sim using a Carter robot within a simulated construction environment. Performance was evaluated using occlusion cost as a safety metric and an efficiency rate based on load, time, and distance.

\subsubsection{Strengths}

The study presents several important strengths. It introduces a novel concept of predicting whether obstacles, particularly human workers, are likely to actively move away, rather than treating all obstacles as passive entities. This intention-aware approach improves both safety and navigation efficiency without requiring modifications to existing planning algorithms, making it practical for integration into current robotic systems.

Additionally, the use of hybrid perception through the combination of camera-based detection and LiDAR sensing enhances environmental understanding and obstacle classification. By predicting obstacle intention, the system reduces unnecessary path replanning when obstacles are expected to clear the path, resulting in smoother robot motion. Experimental results further demonstrate reduced collision risk, measured through occlusion cost, and improved task completion efficiency compared to conventional navigation methods.

\subsubsection{Limitations}

The study's validation was conducted entirely within the NVIDIA Isaac Sim simulation
environment using a Carter robot platform; no physical deployment or real-world testing
was reported, leaving the system's behaviour under real sensor noise, lighting
variation, and hardware latency unverified. The obstacle classification is limited to a
fixed set of construction-site categories (workers, carts, machines, materials), and
the intention prediction concerns only whether an obstacle is likely to \textit{passively}
move out of the way — the framework does not incorporate any explicit, human-initiated
command channel such as a gesture or verbal instruction.

\subsubsection{Research Gap}

Li et al.'s framework demonstrates that feeding predicted intent into an existing Nav2
costmap improves navigation smoothness without modifying the underlying planner — an
architecturally useful insight — but the system has no mechanism for a human to
directly communicate an explicit command to the robot; all "intention" here is inferred
passively from obstacle type and motion, never confirmed through deliberate gesture.
This leaves a gap between passive intention-aware navigation and active human-robot
command interfaces, a gap this project addresses by combining an explicit gesture
recognition channel with the robot's autonomous navigation stack, and by validating the
resulting system on physical hardware rather than in simulation alone.

Across the seven works reviewed, a consistent pattern emerges: systems that achieve
strong intention-prediction performance (Tsitos et al.) tend to operate in isolation
from navigation and obstacle avoidance; systems that incorporate rich, multi-modal
gesture vocabularies (Mahmud et al.) depend on specialized depth-sensing hardware and
remain validated only in simulation; systems that successfully integrate intention
prediction into an existing navigation stack (Li et al.; Yu et al.) do so using
passive, inferred signals rather than explicit human commands, and are likewise
validated only in simulation; and Muhtadin et al., while methodologically the closest
comparison to the present work -- a MediaPipe-based, ROS2-integrated, edge-deployed
gesture pipeline -- validate their system only on a stationary robotic arm, with no
mobile platform, navigation, or SLAM-based mapping involved. No reviewed work combines
an explicit, gesture-based command channel with autonomous navigation on physically
deployed, edge-only hardware.
This project addresses that gap by validating explicit gesture recognition and
physical SLAM-based environment mapping on real hardware, with autonomous
path-planning (Nav2) architecturally integrated into the same pipeline; full
navigation validation is reported as ongoing work in Chapter 4.

\begin{landscape}

\section{Literature Review References}

\footnotesize

\begin{longtable}{|
>{\RaggedRight\arraybackslash}p{3cm}|
>{\RaggedRight\arraybackslash}p{3cm}|
>{\RaggedRight\arraybackslash}p{5cm}|
>{\RaggedRight\arraybackslash}p{4cm}|
>{\RaggedRight\arraybackslash}p{4cm}|
>{\RaggedRight\arraybackslash}p{4cm}|
}
\caption{Comparative Summary of Additional Reviewed Literature}
\label{tab:lit_review_landscape}\\
\hline
\textbf{AUTHORS/YEAR} & 
\textbf{RESEARCH TOPIC} & 
\textbf{METHODOLOGY} & 
\textbf{LIMITATIONS} & 
\textbf{RESEARCH GAP} & 
\textbf{IMPROVEMENT} \\
\hline
\endfirsthead

\hline
\textbf{Authors/Year} & 
\textbf{Topic} & 
\textbf{Methodology} & 
\textbf{Limitations} & 
\textbf{Research Gap} & 
\textbf{Improvement} \\
\hline
\endhead

Xinbo Yu et al. (2025) \cite{yu2024} &
Real-Time Trajectory Planning and Obstacle Avoidance for Human--Robot Co-Transporting &
The study proposes a real-time human–robot co-transporting system that uses a slider-based mechanism and 2D LiDAR with a multiscale perception model to estimate human motion, plan trajectories, and generate safe, smooth navigation with obstacle avoidance. &

The study is limited by its reliance on accurate sensor measurements, use of 2D LiDAR with restricted perception, and indirect estimation of human intention, which may affect robustness in complex environments.&

This paper determines intention through a mechanical slider and rotating disc rather than directly understanding human motion and gestures. &
Future work can integrate multimodal intention recognition combining motion, gesture, and contextual awareness to improve prediction accuracy and adaptability. \\
\hline

Javier Laplaza et al. (2022) \cite{laplaza2022} &
Context and Intention for 3D Human Motion Prediction in Handover Tasks &
The study developed an attention-based deep learning model that predicts human motion and intention using multiple inputs such as motion history, robot position, and environmental context. &
The experiments were limited to handover scenarios and controlled datasets, which may not generalize well to broader real-world interactions. &
The approach is restricted to specific collaborative tasks and does not fully address general human intention recognition across different environments. &
The system can be improved by extending it to diverse real-world scenarios and integrating real-time sensing for more adaptive and generalizable predictions. \\
\hline

Yongshi Liang and Pai Zheng (2024) \cite{liang2024} &
Human Intention Prediction with Visual-Language Model &
The study proposed a visual-language model for predicting human intention using multimodal data from vision and language inputs in structured environments. &
The model requires high computational resources and was mainly tested in traffic and pedestrian scenarios without real-time robotic integration. &
The work focuses on intention prediction but does not incorporate gesture recognition or navigation for human--robot collaboration. &
Future improvements can focus on integrating gesture recognition and real-time robotic control to enable practical human--robot interaction. \\
\hline

Muhtadin et al. (2025) \cite{muhtadin2025} &
Hand Gesture Recognition for Collaborative Robots Using Lightweight Deep Learning in Real-Time Robotic Systems &
The study developed a MediaPipe-based hand landmark extraction pipeline feeding a lightweight neural network classifier, integrated with ROS2 to control a UR5 robotic arm in real time. &
The system is deployed on a stationary robotic arm with no mobile platform, no autonomous navigation, and no SLAM-based environment mapping. &
The work demonstrates a closely related edge-deployed gesture pipeline but does not extend to mobile robotics or navigation integration. &
Future work could extend the lightweight gesture pipeline to mobile platforms with navigation and environment mapping capability. \\
\hline

\end{longtable}

\normalsize

\end{landscape}
\cleardoublepage

\chapter{System Design and Methodology}

\section{Introduction}
This chapter details the methodology, system design, and physical implementation of the real-time, gesture-controlled autonomous mobile robot developed for this study. Moving beyond theoretical simulation, this project is engineered as a fully edge-deployed system. It executes a complex visual perception pipeline natively on onboard embedded hardware, while spatial mapping was validated within a parallel simulation environment, and the autonomous navigation framework was architecturally integrated for future execution.

The chapter provides a comprehensive breakdown of the system components, progressing from the underlying physical hardware and middleware deployment strategies to the high-level cognitive processes. Specifically, it outlines the architecture of the vision-based intention recognition pipeline—detailing spatial person detection, hand landmark extraction, and geometric feature engineering. Furthermore, it describes the centralised decision-making state machine and the architectural integration of Simultaneous Localization and Mapping (SLAM) for autonomous navigation. By detailing these interconnected subsystems, this chapter documents how human hand gestures were reliably captured, classified, and translated into safe robotic motion within this study.

\section{System Requirements}

Before detailing the system architecture, this section specifies the hardware and
software components required to reproduce this system, and the reasoning behind each
selection.

\subsection{Hardware Requirements}

Table 3.1 lists the physical components used in this project. Costs are omitted here
by design decision; component specifications are provided instead to support
replication.

\begin{table}[H]
\centering
\caption{Hardware Components and Specifications}
\label{tab:hardware_requirements}

\begin{tabular}{|p{4cm}|p{10cm}|}
\hline
\textbf{Component} & \textbf{Specification} \\
\hline

Primary computing unit & Raspberry Pi 5, 8GB RAM \\
\hline

Robot expansion board & ESP32-S3 co-processor; 4-channel encoder motor driver; 2-channel PWM servo driver; 1-channel LiDAR interface; 6-axis IMU \\
\hline

Drive motors & 4$\times$310 encoder motors (metal gearbox, Hall-effect encoder) \\
\hline

LiDAR sensor & MS200 Time-of-Flight LiDAR; 12m range; 30 Klux outdoor light resistance \\
\hline

Camera & 2MP USB monocular camera; 2-DOF pan-tilt gimbal mount \\
\hline

Battery & 7.4V 2000mAh rechargeable Li-ion battery \\
\hline

Chassis & Yahboom aluminium-alloy differential-drive platform \\
\hline

\end{tabular}

\end{table}
The Raspberry Pi 5 was initially provisioned with 2GB of RAM; the rationale for the
subsequent upgrade to 8GB is detailed later in this chapter (see Section
\ref{sec:ram_planning}).

\subsection{Software Requirements}

Table 3.2 lists the software stack, distinguishing components that run on the
development workstation from those deployed on the robot itself, reflecting the
architectural split described in Section~\ref{sec:software_framework}.

\begin{table}[H]
\centering
\caption{Software Components and Deployment Target}
\label{tab:software_requirements}

\renewcommand{\arraystretch}{1.2}

\begin{tabular}{|p{3.2cm}|p{3cm}|p{7cm}|}
\hline
\textbf{Component} & \textbf{Deployment Target} & \textbf{Role} \\
\hline

Ubuntu 24.04 & Development workstation & Host OS (Lenovo V15 ADA) \\
\hline

ROS2 Jazzy Jalisco & Development workstation & Development middleware \\
\hline

ROS2 Humble Hawksbill & Robot (Docker) & Deployed middleware, LTS release \\
\hline

Docker Engine & Robot & Container isolation for the three robot containers \\
\hline

MediaPipe HandLandmarker & Robot & Hand landmark extraction \\
\hline

Ultralytics YOLOv8n & Robot & Person and spatial-zone detection \\
\hline

ONNX Runtime & Robot & Gesture classifier and LSTM path predictor inference \\
\hline

scikit-learn (joblib) & Robot & Gesture label decoding \\
\hline

SLAM Toolbox & Robot & Occupancy grid mapping \\
\hline

Navigation2 (Nav2) & Robot & Autonomous path planning \\
\hline

micro-ROS & ESP32-S3 & Firmware and UART bridge to ROS2 graph \\
\hline

Python 3 & Both & Implementation language for all custom ROS2 nodes \\
\hline

\end{tabular}
\end{table}

This split reflects the resource-planning discussion in Section~\ref{sec:ram_planning}: development
convenience on the workstation is decoupled from the stability requirements of the
physically deployed robot, with Docker providing the isolation boundary between the
two ROS2 distributions.

\section{System Architecture Overview}
The system architecture of the autonomous mobile robot is divided into two main components: a high-level processing unit responsible for perception and decision-making, and a low-level hardware controller responsible for driving the motors and reading basic sensors. These hardware components are unified through an isolated software deployment strategy to ensure reliability in real-world environments.

\subsection{Hardware Platform Specification}
The physical foundation of the system is the Yahboom MicroROS-Pi5 differential-drive mobile platform. To achieve the objective of localized, edge-based processing, the primary computational unit is a Raspberry Pi 5. This single-board computer serves as the edge processor for high-level machine learning inference, computer vision tasks, and core robotic middleware management. 

Low-level hardware management and power distribution are handled by a custom robotics expansion board featuring an integrated ESP32-S3 microcontroller. This expansion board acts as a hardware bridge, interfacing directly with the robot's motor drivers, wheel encoders, and an onboard Inertial Measurement Unit (IMU). The microcontroller continuously aggregates this raw odometry and inertial data, transmitting it to the Raspberry Pi to maintain the robot's spatial coordinate frame. For environmental perception, the system utilizes an MS200 2D Light Detection and Ranging (LiDAR) sensor, which provides 360-degree planar distance scans at a variable scan rate. Visual data for human intention recognition is captured via a standard Universal Serial Bus (USB) monocular camera mounted on a two-axis (pan-tilt) servo mechanism, allowing for direct frontal observation of the operator. 

\begin{figure}[H]
\centering
\includegraphics[width=0.6\textwidth]{hardware design.png}
\caption{Hardware Connection Diagram}
\label{fig:hardware_connection}
\end{figure}

\subsection{Software Framework and Deployment Strategy}
\label{sec:software_framework}
The software architecture is built natively on the Robot Operating System 2 (ROS2 Humble Hawksbill distribution) \cite{macenski2022}. ROS2 was selected for its reliable, low-latency Data Distribution Service (DDS) communication between nodes. To optimize performance, the system employs a multi-language approach: high-level cognitive nodes for perception and decision-making are implemented in Python to leverage its extensive machine learning ecosystems, while the low-level hardware bridging relies on the ESP32-S3's pre-configured manufacturer firmware.

Because the ESP32-S3 microcontroller lacks the architecture to support a full Linux-based environment required for standard ROS2, it operates a micro-ROS node. This node transmits raw sensor data over a high-speed serial connection to a micro-ROS agent hosted on the Raspberry Pi, seamlessly linking the microcontroller into the main ROS2 network.

To ensure the software runs reliably and without dependency conflicts, the high-level application stack is packaged using Docker containers. The architecture is deployed across three isolated containers on the Raspberry Pi 5:
\begin{enumerate}
    \item \textbf{\texttt{yahboom\_base}}: Manages core hardware communication, serial interfacing, and LiDAR data publishing.
    \item \textbf{\texttt{micro\_ros\_agent}}: Translates data between the ESP32-S3 microcontroller and the main ROS2 network.
    \item \textbf{\texttt{yahboom\_gesture}}: Houses the custom Python-based perception, spatial filtering, and decision-making nodes developed for this study.
\end{enumerate}

All three containers are built on Yahboom's official \texttt{yahboomtechnology/ros-humble:4.1.2} base image, with the custom perception and decision-making nodes for this study layered into the \texttt{yahboom\_gesture} container on top of that base.

All three containers are launched in Docker's \texttt{host} network mode, meaning they share the Raspberry Pi's network stack directly rather than operating on an isolated container network. This allows ROS2's DDS-based discovery to function across container boundaries exactly as it would between ordinary host processes, so \texttt{yahboom\_base}, \texttt{micro\_ros\_agent}, and \texttt{yahboom\_gesture} locate and communicate with one another over the shared \texttt{ROS\_DOMAIN\_ID} without any additional port-forwarding or bridge configuration.

Finally, to ensure the robot is operational upon power cycling, the pipeline is managed by a background system service (\texttt{cognition.service}). This service ensures that the gesture recognition and core base drivers launch automatically the moment the robot is turned on. While the robot operates headlessly during runtime to preserve processing memory, secure remote access is maintained via Secure Shell (SSH) for administration, and RViz2 is used on a remote developer workstation for live spatial monitoring. 

In alignment with standard robotic development pipelines, environment mapping using SLAM Toolbox was successfully validated through physical deployment, as detailed in Section~\ref{sec:ram_planning}. While the Navigation2 (Nav2) framework was architecturally integrated into the software stack, its physical and simulated execution was constrained by map generation limits, the diagnostics of which are detailed in Chapter 4.

\textit{[Insert Figure Y: System deployment diagram showing the three Docker containers, systemd, RViz2 remote monitoring, and the simulation interface]}

\subsection{Hardware Resource Planning and Memory Constraints}
\label{sec:ram_planning}

The initial hardware configuration for the Raspberry Pi 5 used a 2GB RAM variant,
selected on the basis of cost and the assumption that the perception pipeline's
individual components were lightweight enough to run concurrently. During integration
testing of the SLAM subsystem — specifically, while attempting to build a live
occupancy grid map using SLAM Toolbox while the gesture recognition pipeline (YOLOv8n,
MediaPipe HandLandmarker, and the MLP classifier) remained active — the system began
experiencing repeated \texttt{std::bad\_alloc} failures, manifesting as ROS2 nodes
terminating abruptly (\texttt{Aborted}) and topics such as \texttt{/scan} and
\texttt{/camera/image\_raw} disappearing intermittently. Diagnostic investigation traced
this to Linux out-of-memory (OOM) killer intervention: the combined working set of the
perception stack, SLAM mapping process, and micro-ROS agent exceeded the 2GB board's
available memory under concurrent load, causing the kernel to terminate processes
without warning.

This finding led to a deliberate hardware revision: the Raspberry Pi 5 was upgraded to
the 8GB RAM variant. This resolved the OOM-related crashes and allowed SLAM mapping to
run concurrently with the full gesture recognition pipeline without process
termination. This experience is documented here because it directly informed the
system's final resource budget and is referenced again in the Quality-of-Service design
decisions in Section~\ref{sec:ros2_middleware}, where topic-level memory pressure was further mitigated
through QoS profile selection.

\section{Data Acquisition and Dataset Preparation}
To ensure the Multilayer Perceptron (MLP) gesture classifier performed reliably during physical deployment, a custom dataset was engineered specifically for this study. Utilizing pre-existing public datasets often introduces a ``domain shift'' error, where the training images do not match the focal length or distortion of the robot's actual hardware. To eliminate this, all training data was captured directly through the robot's onboard Universal Serial Bus (USB) monocular camera.

\subsection{Recording Conditions and Class Balancing}
\label{sec:recording_conditions}
The dataset was constructed to encompass the six explicit operational commands: STOP, GO, LEFT, RIGHT, BACK, and FOLLOW. To ensure the neural network learned the geometric relationships of the hand rather than memorizing the background, data collection was conducted across three distinct indoor environments with varying conditions of natural and artificial fluorescent lighting. 

To prevent model bias toward any specific gesture, the dataset was strictly balanced. Exactly 1,000 distinct frames were recorded for each of the six gesture classes, resulting in a total dataset of 6,000 labeled instances. During recording, the operator performed micro-variations of each gesture (altering wrist roll, pitch, and yaw within a 15-degree tolerance) to ensure the model generalized effectively to imperfect human inputs.

\subsection{Preprocessing and Train-Test Splitting}
The raw captured frames were not fed directly into the model. To mirror the live inference pipeline exactly, the raw dataset was processed through the YOLOv8n spatial filter to generate bounding boxes, followed by the MediaPipe HandLandmarker to extract the 21 anatomical coordinates. Frames where MediaPipe failed to confidently extract the hand were discarded and re-recorded. The extracted coordinates were then mathematically converted into the 19 normalized geometric features.

Prior to training, the finalized 6,000-sample feature dataset was randomly shuffled and partitioned using a standard machine learning distribution ratio: 70\% (4,200 samples) allocated for training the model, 15\% (900 samples) reserved for continuous validation during training, and 15\% (900 samples) isolated as a blind test set to evaluate final classification accuracy. 

\subsection{Dataset Statistics}

This section provides a detailed tabular breakdown of the gesture dataset described
in Section~\ref{sec:recording_conditions} and Section 3.4.2 (Preprocessing and
Train-Test Splitting).

\begin{table}[H]
\centering
\caption{Per-Class Sample Distribution}
\label{tab:appendix_class_distribution}
\begin{tabular}{|l|c|c|}
\hline
\textbf{Gesture Class} & \textbf{Samples} & \textbf{Share of Dataset} \\
\hline
STOP & 1,000 & 16.7\% \\
\hline
FOLLOW & 1,000 & 16.7\% \\
\hline
GO & 1,000 & 16.7\% \\
\hline
LEFT & 1,000 & 16.7\% \\
\hline
RIGHT & 1,000 & 16.7\% \\
\hline
BACK & 1,000 & 16.7\% \\
\hline
\textbf{Total} & \textbf{6,000} & \textbf{100\%} \\
\hline
\end{tabular}
\end{table}

\begin{table}[H]
\centering
\caption{Train / Validation / Test Split}
\label{tab:appendix_split}
\begin{tabular}{|l|c|c|}
\hline
\textbf{Split} & \textbf{Samples} & \textbf{Proportion} \\
\hline
Training & 4,200 & 70\% \\
\hline
Validation & 900 & 15\% \\
\hline
Test & 900 & 15\% \\
\hline
\textbf{Total} & \textbf{6,000} & \textbf{100\%} \\
\hline
\end{tabular}
\end{table}

Data collection spanned three distinct indoor environments with varying natural and
artificial fluorescent lighting conditions, with micro-variations in wrist roll,
pitch, and yaw (within a 15-degree tolerance) introduced during recording to improve
generalization to imperfect human inputs, as described in
Section~\ref{sec:recording_conditions}.

\section{Perception and Intention Recognition Pipeline}
The perception pipeline was designed as a sequential filter to isolate and classify human hand gestures from raw visual data. This architecture was specifically optimized for localized inference, allowing the system to extract precise features and execute classification models natively on the edge processor. The pipeline isolated a valid human operator, extracted anatomical hand landmarks, engineered scale-invariant geometric features, and classified the intended gesture.

\subsection{Spatial Zone Person Detection}
The first stage of the pipeline utilized the YOLOv8n (nano variant) object detection model \cite{yolov8_ultralytics}. During initial development and empirical testing, a critical vulnerability was observed: the downstream skeletal tracking framework would unpredictably extract landmarks from any visible hand in the camera frame. In environments with multiple individuals, the robot would erroneously lock onto background movements or gestures performed by non-operators, leading to erratic system behavior.

To resolve this interference and enforce operational safety, a spatial zone filter was engineered over the raw YOLOv8n outputs. The camera's field of view was programmatically segmented into an active acceptance zone spanning 45\% of the frame width and 65\% of the frame height, centred on the frame origin. The system evaluated the center coordinates of all detected ``person'' bounding boxes; any detection falling outside this boundary was explicitly ignored. To support scenarios with multiple bystanders near the operational zone, the system
tracks up to three candidate persons simultaneously for motion-history purposes, but
selects a single active operator per frame as the largest (closest) zone-eligible
detection. Once this operator produces a validated gesture command (per the 5-frame
majority-vote buffer described in Section~\ref{sec:temporal_buffering}),
\texttt{brain\_node} grants a 3-second lock on that subject, during which newly
appearing candidates are ignored regardless of size or proximity. This logical filtering ensured the robot interacted exclusively with an operator positioned directly in its operational path. 

It is acknowledged that this geometric heuristic relied on spatial positioning rather than unique operator identification. In edge cases where multiple individuals occupied the active acceptance zone simultaneously, system disambiguation could fail. Expanding this architecture to include facial recognition or persistent biometric identity tracking remains an area for future development. A standalone facial recognition prototype was evaluated offline using InsightFace (buffalo\_sc, \texttt{w600k\_mbf.onnx}), successfully distinguishing between two enrolled operator identities with confidence scores of 0.54--0.57 against an ``Unknown'' baseline of 0.09. This prototype was not integrated into the deployed pipeline and remains a candidate for future work in resolving multi-operator disambiguation.

\textit{[Insert Figure Z: Spatial zone filter diagram representing the camera frame with the acceptance zone highlighted]}

\subsection{Vision-Based Hand Tracking}
Once an operator was successfully isolated within the spatial zone, the pipeline engaged the MediaPipe HandLandmarker framework \cite{lugaresi2019mediapipe} for skeletal tracking. This framework extracted high-fidelity anatomical data from the 2D monocular camera feed.

For each valid frame containing the primary operator, the node identified the spatial coordinates of 21 specific hand landmarks, including the wrist, knuckles, and fingertips. Each landmark was output as a normalized three-dimensional coordinate (x, y, z). While this provided a detailed topological representation of the hand, feeding these raw, unscaled coordinates directly into a lightweight classifier introduced significant statistical variance depending on the operator's physical distance from the camera. Consequently, these raw landmarks served strictly as the foundational input for the subsequent geometric feature engineering stage.

\subsection{Geometric Feature Engineering}
To ensure the classification model remained robust regardless of the operator's distance from the camera, the raw spatial landmarks were transformed into a set of scale-invariant geometric features. 

Instead of utilizing absolute coordinates, the system engineered 19 specific geometric relationships based on the relative distances and angles between critical hand joints. The wrist landmark (Landmark 0) was established as the primary structural origin point. From this origin, the Euclidean distance to the terminal fingertips was calculated using the standard 3D distance formula:

\begin{equation}
D = \sqrt{(x_2 - x_1)^2 + (y_2 - y_1)^2 + (z_2 - z_1)^2}
\end{equation}

To achieve true scale invariance, these calculated distances were mathematically normalized against the primary bounding box of the hand itself. This normalization ensured that a gesture performed one meter away from the camera generated the identical 19-value numerical array as a gesture performed three meters away. This engineered 1D array of 19 optimized features served as the direct input for the gesture classification model.

\begin{table}[H]
\centering
\caption{Engineered Geometric Features (19 total)}
\label{tab:geometric_features}
\begin{tabular}{|c|l|l|}
\hline
\textbf{\#} & \textbf{Feature} & \textbf{Description} \\
\hline
1 & Thumb curl angle & Angle between MCP--PIP and PIP--TIP segments \\
\hline
2 & Index curl angle & Angle between MCP--PIP and PIP--TIP segments \\
\hline
3 & Middle curl angle & Angle between MCP--PIP and PIP--TIP segments \\
\hline
4 & Ring curl angle & Angle between MCP--PIP and PIP--TIP segments \\
\hline
5 & Pinky curl angle & Angle between MCP--PIP and PIP--TIP segments \\
\hline
6 & Thumb--Index spread & Angle between thumb and index tip vectors (from wrist) \\
\hline
7 & Index--Middle spread & Angle between index and middle tip vectors (from wrist) \\
\hline
8 & Middle--Ring spread & Angle between middle and ring tip vectors (from wrist) \\
\hline
9 & Thumb tip distance & Thumb tip--to--wrist distance, normalized by palm width \\
\hline
10 & Index tip distance & Index tip--to--wrist distance, normalized by palm width \\
\hline
11 & Middle tip distance & Middle tip--to--wrist distance, normalized by palm width \\
\hline
12 & Ring tip distance & Ring tip--to--wrist distance, normalized by palm width \\
\hline
13 & Pinky tip distance & Pinky tip--to--wrist distance, normalized by palm width \\
\hline
14 & Palm angle & Angle of middle-MCP--to--wrist vector vs. vertical reference \\
\hline
15 & Index pointing angle & 2D angle of the index finger's direction vector \\
\hline
16 & Thumb pointing angle & 2D angle of the thumb's direction vector \\
\hline
17 & Index--Thumb angle & Angle between index and thumb vectors (from wrist) \\
\hline
18 & Thumb--Index tip distance & Distance between thumb tip and index tip, normalized \\
\hline
19 & Thumb height & Thumb tip depth (z-axis) relative to wrist \\
\hline
\end{tabular}
\end{table}

Palm width (used for normalization in features 9--13 and 18) is computed as the Euclidean distance between the index and pinky MCP landmarks.

\subsection{Gesture Classification Model}
\label{sec:gesture_classification_model}
An initial implementation of the gesture classification stage used a rule-based classifier: fixed geometric thresholds were applied directly to the 19 engineered features (e.g., "if thumb-to-index distance exceeds threshold X, classify as STOP"). While functional for a small number of clearly distinguishable gestures, this approach proved brittle in practice — small variations in hand orientation, distance from camera, and inter-operator differences in gesture execution frequently pushed feature values across hand-tuned thresholds, producing inconsistent classifications that required constant manual recalibration. This limitation motivated the transition to a trained Multilayer Perceptron (MLP) classifier, which could learn decision boundaries directly from labelled examples rather than relying on manually specified rules, and which generalized more reliably across the natural variation present in human gesture execution.

With the spatial data reduced to 19 scale-invariant features, the final stage of the perception pipeline was the classification of the intentional command using this MLP. The model was configured to classify the incoming feature array into one of six discrete operational commands: STOP, GO, LEFT, RIGHT, BACK, and FOLLOW.

Two rounds of dataset collection and training were conducted. An initial dataset of 3,000 samples (unbalanced across classes) produced a classifier with 99.67\% test accuracy; however, to improve robustness against the operator-variation issues that motivated the move away from the rule-based approach in the first place, the dataset was subsequently expanded to a balanced 6,000-sample set (1,000 samples per class, described in Section~\ref{sec:recording_conditions}), yielding a final test accuracy of 99.38\% — a marginal reduction attributable to the harder, more balanced classification task rather than a regression in model quality.

The final trained classifier was exported to Open Neural Network Exchange (ONNX)
format (\texttt{gesture\_model.onnx}) and loaded via ONNX Runtime
(\texttt{onnxruntime.InferenceSession}) inside \texttt{gesture\_node} at runtime, with
a separate \texttt{label\_encoder\_pi.pkl} file -- loaded via scikit-learn's joblib --
handling decoding of the model's integer output back to a gesture label string. ONNX
Runtime was used consistently across both learned models in the perception pipeline,
providing a uniform inference approach for the gesture classifier and the LSTM path
predictor alike. The raw classification output was then published to the ROS2 middleware, where it was intercepted by the centralized decision-making node.

\textit{[Insert Figure: Feature extraction diagram illustrating the conversion of the 21 MediaPipe landmarks into the 19 features, feeding into the MLP classifier]}

\subsection{Human Movement Prediction (LSTM Path Predictor)}
\label{sec:lstm_predictor}
In addition to static hand-gesture classification, a secondary recognition pathway was implemented to capture human movement intent over time, complementing the frame-by-frame gesture pipeline with a temporal model of operator motion. Sequential pose landmarks extracted via MediaPipe Pose were fed into a Long Short-Term Memory (LSTM) network trained to predict an operator's short-horizon future position, allowing the system to anticipate movement direction rather than react only after a position change had already occurred. The model was trained and evaluated on 152 motion sequences, achieving an Average Displacement Error (ADE) of 25.8 pixels and a Final Displacement Error (FDE) of 32.4 pixels on held-out test sequences. The trained model was exported to ONNX format, consistent with the deployment strategy used for the gesture classifier, to ensure uniform inference performance on the Raspberry Pi 5's edge processor.

\begin{figure}[H]
\centering
\includegraphics[width=0.85\textwidth]{path_predictor_loss.png}
\caption{LSTM Path Predictor Training Loss}
\label{fig:lstm_loss}
\end{figure}

\begin{figure}[H]
\centering
\includegraphics[width=0.85\textwidth]{path_prediction_examples.png}
\caption{LSTM Path Prediction — History vs. True vs. Predicted (Sample Sequences)}
\label{fig:lstm_examples}
\end{figure}

\section{Centralized Decision-Making and Control}
The translation of raw classification data into safe, physical robot movement was managed by a centralized Python-based control node (\texttt{brain\_node}). This node functioned as the core state machine, evaluating the continuous stream of AI predictions and deciding when to issue target velocities to the motor controller via the \texttt{/cmd\_vel} ROS2 topic.

\subsection{The Brain Node Architecture}
The \texttt{brain\_node} operated asynchronously from the perception pipeline. It subscribed to the gesture classification topic and the YOLOv8n spatial bounding box topic. Rather than reacting to every single frame—which would result in erratic jitter—the node processed these inputs through strict logical conditions using a timer-based control loop (\texttt{create\_timer}), adopted over a manual polling loop to guarantee reliable, consistent callback execution. It maintained discrete operational states (IDLE, WAITING\_CONFIRMATION, LOCKED, and EXECUTING) to ensure the robot transitioned smoothly between stationary observation and active motion.

\subsection{Temporal Buffering and Locking Logic}
\label{sec:temporal_buffering}
Machine learning models operating on live video feeds occasionally produce
single-frame misclassifications due to motion blur or lighting shifts. To prevent the
robot from responding to a false positive, two independent safeguards were
implemented. First, \texttt{gesture\_node} rejects any single-frame prediction with a
confidence score below 0.65 before it is published. Second, \texttt{brain\_node}
maintains a 5-frame rolling buffer of recently published gesture IDs and confirms a
gesture once it accounts for at least 3 of the 5 most recent buffered predictions -- a
majority-vote confirmation scheme rather than a requirement for five identical
consecutive frames.

Once a command was validated, the node initiated a 3-second subject lock. During this period, the system locked onto the current operator and ignored any conflicting gestures or new persons entering the frame. This locking mechanism prevented bystanders from accidentally interrupting an active command sequence and enforced predictable human-robot interaction.

In addition to this gesture-driven state machine, \texttt{brain\_node} exposes a custom \texttt{emergency\_stop} ROS2 service, allowing an operator or external supervisory process to immediately override all classification-driven behavior and command a zero-velocity state regardless of the current temporal buffer or lock status.

\subsection{Reactive Person Tracking (Visual Servoing)}
When the operator issued the ``FOLLOW'' command, the robot initiated a reactive person-tracking sequence. This behavior did not utilize predictive pathing algorithms; rather, it employed visual servoing based on proportional control mathematics. 

The system tracked the center coordinates of the operator's YOLOv8n bounding box ($B_x$) and compared them to the fixed center of the camera frame ($C_x$). The error between these two points dictated the required angular velocity ($v_\omega$) to keep the human centered in the field of view, calculated using a proportional tuning constant ($K_p$):

\begin{equation}
v_\omega = K_p \times (C_x - B_x)
\end{equation}

Simultaneously, the area of the bounding box was monitored to estimate proximity. If the bounding box area dropped below a predetermined threshold (indicating the operator was moving further away), a positive linear velocity was commanded. If the bounding box area exceeded the threshold (indicating the operator was too close), the robot stopped to maintain a safe operating distance. 

\textit{[Insert Figure: UML Brain node state machine diagram showing states and transitions]}

\section{ROS2 Middleware Integration and Topic Interfaces}
\label{sec:ros2_middleware}
The distributed nodes of the perception, decision-making, and hardware abstraction layers were bound together by the ROS2 Data Distribution Service (DDS). To ensure real-time performance without overwhelming the Raspberry Pi 5's available RAM, specific Quality of Service (QoS) profiles and message typologies were engineered for the communication topics.

\subsection{ROS2 Domain Isolation}
Early integration testing revealed a subtle configuration fault in which the base hardware driver package (\texttt{yahboom\_base}) operated on \texttt{ROS\_DOMAIN\_ID=0} while the cognition package (\texttt{yahboom\_gesture}) operated on \texttt{ROS\_DOMAIN\_ID=20}. Because ROS2 nodes on different domain IDs cannot discover or communicate with one another by design, this mismatch silently prevented \texttt{brain\_node} from receiving expected control inputs, resulting in a default zero-velocity command being published at 10Hz regardless of gesture state. Resolving this required explicitly aligning both packages to a single shared domain ID, after which command propagation to \texttt{/cmd\_vel} functioned as intended. This is documented here as it directly shaped the final middleware configuration: all cognition and hardware-abstraction nodes are now launched with an explicit, shared \texttt{ROS\_DOMAIN\_ID} set at the container level rather than relying on environment defaults.

\subsection{Sensor Data Streaming}
High-bandwidth environmental data requires low latency. The monocular camera node published visual data to the \texttt{/camera/image\_raw/compressed} topic at a constrained rate of 20 Hertz (Hz) using the standard \texttt{sensor\_msgs/Image} message type. Simultaneously, the MS200 LiDAR published spatial distance arrays to the \texttt{/scan} topic at 12Hz utilizing the \texttt{sensor\_msgs/LaserScan} format.

To prevent network congestion, these continuous sensor streams utilized a ``Best Effort'' QoS profile. This configuration instructed the middleware to drop outdated frames rather than queuing them if processing fell behind, ensuring the AI models always evaluated the most current physical reality rather than processing a backlog of old images.

\subsection{Cognitive and Control Topologies}
Unlike high-frequency sensor streams, cognitive outputs and motor controls require absolute transmission reliability. When the MLP classifier processed a frame, it published a custom \texttt{cognition\_interfaces/Gesture} message, containing an integer gesture ID, a string label, a floating-point confidence score, and a timestamp — to the \texttt{/cognition/gesture} topic at approximately 10Hz. 

Similarly, \texttt{person\_detection\_node} published a custom \texttt{cognition\_interfaces/Detection} message to the \texttt{/cognition/detection} topic, containing the operator's bounding-box center coordinates, a kinematic motion-direction classification, and — following the LSTM integration described in Section~\ref{sec:lstm_predictor} — the model's predicted future x-coordinate (\texttt{predicted\_x}), all at the node's inference rate.

The \texttt{brain\_node} intercepted this data, executed its temporal buffering and state machine logic, and subsequently published physical movement instructions to the \texttt{/cmd\_vel} topic. These motor commands utilized the standard \texttt{geometry\_msgs/Twist} message type, which specified linear velocity on the x-axis and angular velocity on the z-axis. Because a dropped motor command could result in a physical collision, both the cognitive and control topics were configured with a ``Reliable'' QoS profile, guaranteeing message delivery to the micro-ROS agent bridging the ESP32-S3 hardware controller.

\section{Mapping and Autonomous Navigation}
In addition to direct gesture control, the software architecture incorporated standard ROS2 frameworks to map the operational environment and enable autonomous pathfinding.

\subsection{SLAM Toolbox Configuration}
Simultaneous Localization and Mapping (SLAM) was integrated using the ROS2 SLAM Toolbox. The system was configured to ingest 360-degree laser scan data from the MS200 LiDAR alongside wheel odometry aggregated by the ESP32-S3 microcontroller. Using these inputs, the SLAM node generated a 2D occupancy grid map of the surrounding environment, assigning a probability of an obstacle to every 5cm spatial cell. This mapping capability was successfully validated through physical deployment, producing a coherent occupancy grid map of a real indoor room following the RAM upgrade described in Section~\ref{sec:ram_planning}. During the mapping process itself, the robot was manually driven around the target
environment using a joystick teleoperation node (\texttt{joy\_node}/\texttt{joy\_ctrl}),
publishing directly to \texttt{/cmd\_vel} to traverse the space while SLAM Toolbox
constructed the occupancy grid from the resulting LiDAR scans.

\paragraph{SLAM reliability diagnosis and correction.}
\label{sec:slam_reliability_fixes}
Following the initial mapping validation, a subsequent diagnostic session
identified two distinct root causes affecting SLAM output quality, each
investigated and resolved independently.

First, a captured mapping run showed a complete (100\%) drop rate of laser scans
within \texttt{slam\_toolbox}'s message filter. Investigation traced this to a
transform race condition: a static transform publisher was broadcasting the
\texttt{odom\_frame} $\to$ \texttt{base\_footprint} transform at a fixed identity
value simultaneously with the Extended Kalman Filter (EKF) publishing the same
transform dynamically, causing transform lookups to become unreliable. Removing
the conflicting static publisher and designating the EKF as the sole broadcaster
of this transform reduced the scan drop rate from 100\% to a range of 12--80\%
across repeated verification runs.

Second, two independent mapping attempts following this fix both exhibited a
consistent double-lobed, starburst-artifact map shape -- indicative of systematic
rotational drift rather than a driving-technique inconsistency, since the same
pattern recurred across separate drives. Diagnosis found that the EKF's odometry
configuration had yaw fusion disabled, and the IMU's absolute orientation output
was likewise unused, leaving heading estimation dependent solely on integrated
gyroscopic rate -- a configuration prone to dead-reckoning drift, particularly
during turns. Before applying a fix, this was corroborated with direct evidence
rather than assumed: a live capture of \texttt{/odom\_raw} through an approximately
720-degree physical rotation traced a smooth, correctly-wrapping path, confirming
the underlying sensor data was trustworthy and that the drift originated in the
fusion configuration rather than the sensors themselves. Enabling yaw fusion in
the EKF's odometry configuration resolved the issue; the corrected system traced
approximately three full physical rotations with proper wraparound in subsequent
verification.

A mapping drive conducted after both corrections produced a single, coherent room
shape free of the drift artifacts observed previously, superseding the earlier
map referenced above as the system's current representative mapping result.
Autonomous path planning through this corrected map via Nav2 had not yet been
validated at the time of writing, as detailed in the following subsection and in
Chapter 4.

\subsection{Navigation2 (Nav2) Integration}
To enable autonomous movement to predetermined waypoints, the Navigation2 (Nav2) framework \cite{macenski2020marathon2} was architecturally
integrated into the deployment stack. The framework utilized the Adaptive Monte Carlo Localization (AMCL) algorithm to determine the robot's coordinates on the pre-generated SLAM map, while relying on costmaps to dynamically plan paths around detected obstacles. While the integration of these nodes was successful at the software level, physical and simulated execution of the full Nav2 stack was constrained by memory allocation limits encountered during map generation. The specific diagnostic analysis of these hardware limits is documented in Chapter 4.

\section{ROS2 Node Graph Summary}

Table 3.3 consolidates the ROS2 graph relevant to this study's custom cognition software
and the manual teleoperation subsystem used during SLAM mapping (Section
\ref{sec:ram_planning}), as verified on the deployed system. Vendor-default utility nodes
unrelated to this study's contributions (e.g. IP-address publishing) are omitted.

\begin{table}[H]
\centering
\caption{Deployed ROS2 Nodes, Topics, and Services}
\label{tab:node_graph}
\resizebox{\textwidth}{!}{%
\begin{tabular}{|l|l|}
\hline
\textbf{Category} & \textbf{Items} \\
\hline
Custom Cognition Nodes & \texttt{brain\_node}, \texttt{camera\_pub}, \texttt{gesture\_node}, \texttt{person\_detection\_node} \\
\hline
Teleoperation Nodes & \texttt{joy\_node}, \texttt{joy\_ctrl} (manual driving during SLAM mapping) \\
\hline
Custom Messages & \texttt{cognition\_interfaces/Gesture}, \texttt{cognition\_interfaces/Detection} \\
\hline
Custom Service & \texttt{/brain\_node/emergency\_stop} \\
\hline
Cognition Topics & \texttt{/cognition/gesture}, \texttt{/cognition/detection} \\
\hline
Sensor/Control Topics & \texttt{/camera/image\_raw/compressed}, \texttt{/scan}, \texttt{/cmd\_vel}, \texttt{/imu}, \texttt{/odom\_raw} \\
\hline
Teleoperation Topics & \texttt{/joy}, \texttt{/JoyState}, \texttt{/joy/set\_feedback} \\
\hline
Auxiliary Topics & \texttt{/battery}, \texttt{/beep}, \texttt{/servo\_s1}, \texttt{/servo\_s2} \\
\hline
\end{tabular}
}
\end{table}

All custom message and service definitions are provided by the \texttt{cognition\_interfaces}
package within the \texttt{cognition\_ws} workspace, alongside the \texttt{cognition\_perception}
and \texttt{cognition\_brain} packages housing the perception and decision-making nodes
respectively.

\section{Testing Instrumentation and Data Logging}
\label{sec:testing_instrumentation}
To support the quantitative evaluation presented in Chapter 4, a dedicated logging node was developed to run alongside the deployed pipeline during test sessions. Rather than relying on manual observation and hand-tallied results, this node subscribed to the same \texttt{/cognition/gesture} and \texttt{/cmd\_vel} topics consumed by \texttt{brain\_node}, and recorded each trial — the tester-specified ground-truth gesture, the model's predicted label and confidence score, the resulting linear and angular velocity commands, and inter-prediction latency — to a timestamped CSV file. This logging approach was adopted specifically to enable post-hoc statistical analysis (accuracy, precision, recall, F1-score, and confusion matrix construction, presented in Chapter 4) rather than relying on live pass/fail judgment during testing, which would not have supported rigorous quantitative reporting.

\chapter{TESTING, RESULTS AND ANALYSIS}

\section{Overview}
This chapter presents the testing procedures, results, and performance evaluation of the proposed human intention recognition system. The system was tested to assess its ability to recognize hand gestures, classify human movements, generate appropriate robot responses, and support safe navigation within the operating environment. The performance of the integrated hardware and software components was also evaluated to determine the effectiveness of the overall system.
\section{Hardware Testing}

\subsection{Camera Testing}

The USB monocular camera was tested to verify its ability to capture real-time video streams for gesture and movement recognition. The camera successfully captured human gestures and body movements under normal indoor lighting conditions. However, its performance was affected under poor lighting conditions, where reduced image quality occasionally resulted in inaccurate landmark detection and decreased recognition accuracy. This limitation highlights the dependence of the vision-based system on adequate illumination for reliable operation.

\subsection{Lidar Testing}

The Lidar sensor was tested to evaluate its ability to detect obstacles and provide environmental information for mapping and navigation. During testing, the sensor continuously scanned the surrounding environment and generated distance measurements that were used to identify nearby obstacles. The acquired data supported the robot's perception of its environment and contributed to mapping, localization, and navigation processes. The sensor demonstrated reliable obstacle detection within the testing area, enabling the robot to maintain environmental awareness and support safe navigation.

\subsection{Raspberry Pi 5 Testing}

The Raspberry Pi 5 was tested to evaluate its ability to process camera input, execute machine learning models, and communicate with the robot control subsystem through ROS2. During testing, the Raspberry Pi successfully handled image acquisition, landmark extraction, gesture recognition, movement recognition, and robot-control operations. The device was able to execute the required software components simultaneously while maintaining stable communication with the connected hardware modules. The results demonstrated that the Raspberry Pi 5 provided sufficient computational capability for real-time processing and system integration within the proposed human intention recognition framework.

\section{Gesture Recognition Testing}

\subsection{Test Procedure}

Three participants were asked to perform each gesture ten times in front of the camera. The system's predictions were recorded and compared with the expected gesture labels.

\begin{table}[H]
\centering
\caption{Gesture Recognition Test Results (3 participants, 10 trials each per gesture)}
\label{tab:gesture_results}
\begin{tabular}{|l|c|c|}
\hline
\textbf{Gesture} & \textbf{Trials} & \textbf{Accuracy Range (\%)} \\
\hline
Stop & 30 & 83\% -- 100\% \\
\hline
Follow & 30 & 93\% -- 100\% \\
\hline
Go & 30 & 88\% -- 100\% \\
\hline
Left & 30 & 100\% \\
\hline
Right & 30 & 98\% -- 100\% \\
\hline
Back & 30 & 99\% -- 100\% \\
\hline
\end{tabular}
\end{table}

Three participants performed each of the six gesture classes ten times, for 30 trials
per gesture (n = 3 $\times$ 10). The reported range reflects the spread of per-participant
accuracy across the three testers. Stop showed the widest spread (83--100\%); the
remaining five gestures each stayed within a narrower 7--12 percentage-point band.

\section{Movement Recognition Testing}

\subsection{Test Procedure}

Participants performed approaching, moving away, moving left, moving right, and stationary actions while the system monitored body pose sequences.

\subsection{Results}

\begin{table}[H]
\centering
\caption{Movement Recognition Test Results}
\label{tab:movement_results}
\begin{tabular}{|l|c|}
\hline
\textbf{Movement Class} & \textbf{Result} \\
\hline
Approaching & Successfully Recognized \\
\hline
Moving Away & Successfully Recognized \\
\hline
Moving Left & Successfully Recognized \\
\hline
Moving Right & Successfully Recognized \\
\hline
Stationary & Successfully Recognized \\
\hline
\end{tabular}
\end{table}

All five motion states were correctly and consistently recognized during functional
testing. This qualitative result reflects the level of testing rigor at the time these
trials were conducted; a quantitative breakdown comparable to Table 4.1 requires the
structured per-trial logging described in Section~\ref{sec:testing_instrumentation}.

\section{Robot Response Testing}

\subsection{Gesture-Based Robot Control}

\begin{table}[H]
\centering
\caption{Gesture-Based Robot Control Test Results}
\label{tab:robot_response}
\begin{tabular}{|l|l|l|l|}
\hline
\textbf{Gesture} & \textbf{Expected Response} & \textbf{Actual Response} & \textbf{Status} \\
\hline
Stop & -- & -- & -- \\
\hline
Go & -- & -- & -- \\
\hline
Left & -- & -- & -- \\
\hline
Right & -- & -- & -- \\
\hline
Back & -- & -- & -- \\
\hline
Follow & -- & -- & -- \\
\hline
\end{tabular}
\end{table}
\section{Mapping and Navigation Testing}

\subsection{Environment Mapping}

The robot was deployed in a controlled indoor environment where the LiDAR sensor generated an occupancy grid map of the surrounding area.

\begin{table}[H]
\centering
\caption{Obstacle Avoidance Test Results}
\label{tab:obstacle_avoidance}
\begin{tabular}{|l|l|l|}
\hline
\textbf{Scenario} & \textbf{Expected Outcome} & \textbf{Actual Outcome} \\
\hline
Static Obstacle & -- & -- \\
\hline
Narrow Passage & -- & -- \\
\hline
Multiple Obstacles & -- & -- \\
\hline
\end{tabular}
\end{table}
\section{Performance Evaluation and Benchmarking}

\subsection{Performance Metrics}

This section presents the performance metrics used to evaluate the effectiveness of the proposed human intention recognition system, drawn from both training-time evaluation and live system measurement.

\textbf{Classification performance.} 
The final MLP gesture classifier, trained on the
balanced 6,000-sample dataset described in Section~\ref{sec:recording_conditions},
achieved an overall accuracy of 99.38\%. Table~\ref{tab:precision_recall} reports the
exact per-class precision and recall, both of which stayed at or above 0.987 for
every gesture class, with the confusion matrix shown in
Figure~\ref{fig:confusion_matrix}. This training-time evaluation is distinct from,
and generally consistent with, the live physical-trial accuracy ranges reported in
Table 4.1.

\begin{table}[H]
\centering
\caption{Per-Class Precision and Recall}
\label{tab:precision_recall}
\begin{tabular}{|l|c|c|}
\hline
\textbf{Gesture} & \textbf{Precision} & \textbf{Recall} \\
\hline
BACK   & 0.997 & 0.996 \\
\hline
FOLLOW & 0.999 & 0.987 \\
\hline
GO     & 0.993 & 0.994 \\
\hline
LEFT   & 0.991 & 0.991 \\
\hline
RIGHT  & 0.995 & 0.998 \\
\hline
STOP   & 0.988 & 0.997 \\
\hline
\end{tabular}
\end{table}

\begin{figure}[H]
\centering
\includegraphics[width=0.75\textwidth]{MLP_Confusion_Matrix.png}\caption{Confusion Matrix — MLP Gesture Classifier}
\label{fig:confusion_matrix}
\end{figure}

\begin{table}[H]
\centering
\caption{Real-Time System Throughput (Measured)}
\label{tab:throughput}
\begin{tabular}{|l|c|}
\hline
\textbf{Pipeline Stage} & \textbf{Measured Rate} \\
\hline
Camera capture (\texttt{/camera/image\_raw/compressed}) & 20 Hz \\
\hline
Person detection (YOLOv8n) & 2.7 Hz \\
\hline
Gesture classification (\texttt{/cognition/gesture}) & 10 Hz \\
\hline
LiDAR scan (\texttt{/scan}) & 12.5 Hz \\
\hline
\end{tabular}
\end{table}

The person-detection stage's measured throughput (2.7 Hz) is substantially lower than the camera's native capture rate (20 Hz), reflecting the relative computational cost of YOLOv8n inference compared to the lighter-weight MediaPipe and MLP stages running in parallel. This gap does not affect gesture-command responsiveness, since gesture classification (10 Hz) operates as an independent pipeline stage, but it does bound how quickly the person-detection and LSTM movement-prediction pathway can react to a change in operator position.

\textbf{Hardware resource utilisation.} During full-pipeline operation (all four custom nodes active simultaneously with SLAM Toolbox), the Raspberry Pi 5 reported approximately 311\% aggregate CPU utilisation across its four cores (an average of roughly 78\% per core), reflecting the combined load of YOLOv8n inference, MediaPipe HandLandmarker, the MLP and LSTM models, and SLAM occupancy-grid processing running concurrently. This measurement was taken after the RAM upgrade described in Section~\ref{sec:ram_planning}; the same workload was not sustainable under the original 2GB configuration, as detailed in that section.

\subsection{Benchmarking Analysis}

This section compares the performance of the proposed system against the two most directly comparable reviewed works from Chapter 2 that report quantitative gesture-classification accuracy.

\begin{table}[H]
\centering
\caption{Benchmarking Against Reviewed Gesture Recognition Systems}
\label{tab:benchmarking}
\resizebox{\textwidth}{!}{%
\begin{tabular}{|l|l|c|l|}
\hline
\textbf{System} & \textbf{Classifier} & \textbf{Accuracy} & \textbf{Deployment} \\
\hline
This project & MLP (19 geometric features) & 99.38\% & Physical, edge-only (Raspberry Pi 5) \\
\hline
Mahmud et al. (2022) \cite{mahmud2022} & HMM (best-performing of 3 classifiers) & 94.61\% average, 95.67\% maximum & Simulation only (Kinect 3D joints) \\
\hline
Tsitos et al. (2022) \cite{tsitos2022} & Logistic Regression / SVM / Decision Tree & Not independently verified -- see note & Physical (UR3 arm, reaching task) \\
\hline
\end{tabular}
}
\end{table}

This project's classifier reports higher raw accuracy than Mahmud et al.'s best-performing classifier (94.61\% average across three age categories, 95.67\% maximum for their strongest demographic group, using HMM on 3,600 Kinect-derived samples from 24 participants). However, this comparison carries an important caveat beyond the numbers themselves: Mahmud et al.'s evaluation was conducted entirely in simulation with a depth sensor, while this project's figure comes from a physically deployed, depth-sensor-free system. The more meaningful distinction is therefore architectural rather than purely numerical: this project is validated through live, physical-hardware trials (Table 4.1) on a fully edge-deployed platform, whereas Mahmud et al.'s figure is a training/cross-validation result never deployed to real hardware.

Tsitos et al.'s exact reported accuracy or F1-score could not be independently verified for this comparison and is therefore omitted from the numerical benchmarking above rather than estimated. Their study is retained in the qualitative comparison (Table~\ref{tab:lit_review_landscape}) on the basis of its distinct evaluation approach (a competitive reaching-game task rather than discrete gesture classification), which is not a like-for-like accuracy comparison in any case.

\section{Overall System Evaluation}

Evaluated as a whole, the system demonstrates that explicit gesture recognition and
implicit movement prediction can be combined on a single edge-deployed platform
without reliance on cloud processing or GPU acceleration. The gesture recognition
subsystem's 99.38\% test accuracy translated to real-world physical trial accuracy of
83--100\% across three participants, with five of six gesture classes staying within
a narrow 7--12 percentage-point band -- indicating that laboratory-level classifier
performance transferred reasonably well to physical deployment, with Stop being the
notable exception showing the widest variability.

The movement prediction subsystem's LSTM-based path predictor, achieving an Average
Displacement Error of 25.8px and Final Displacement Error of 32.4px, operated
successfully alongside the gesture pipeline in real time, with its output embedded
directly in the \texttt{Detection} message's \texttt{predicted\_x} field rather than
published as a separate topic. All five qualitative movement classes were
consistently recognized during functional testing.

System-level throughput measurements confirmed the pipeline sustained concurrent
operation of all major subsystems on the Raspberry Pi 5's central processing unit:
camera streaming at 20Hz, person detection at 2.7Hz, gesture classification at
approximately 10Hz, and LiDAR scanning at 12.5Hz, with aggregate CPU utilization of
approximately 311\% across four cores (roughly 78\% average per core). This confirms
the system operates with meaningful computational headroom rather than at saturation,
though the margin is not large given the additional load that would be introduced by
deploying the currently-offline facial recognition prototype.

Benchmarked against comparable systems from the literature review, this project's
gesture classifier reported higher raw accuracy (99.38\%) than Mahmud et
al.'s~\cite{mahmud2022} best-performing classifier (94.61\% average, 95.67\% maximum,
using HMM), though this comparison carries an important caveat: Mahmud et al.'s
evaluation was conducted entirely in simulation using a depth sensor, whereas this
project's figure reflects a physically deployed, depth-sensor-free system operating
under real-world lighting and occlusion conditions. Tsitos et al.'s~\cite{tsitos2022}
exact reported accuracy could not be independently verified and is therefore
excluded from numerical comparison rather than estimated; their system remains a
relevant qualitative comparison point as a physically deployed intention-prediction
system for a reaching task rather than discrete gesture classification.

Taken together, the evaluation supports the conclusion that the system meets its core
functional objectives -- reliable gesture recognition, functioning movement
prediction, and real-time operation on constrained edge hardware -- while also making
clear, through the benchmarking caveats and the still-unvalidated autonomous
navigation component, that "meets objectives" and "fully production-ready" are
distinct claims; the latter is not made here.

\section{Summary}

This chapter presented the empirical results of the intention recognition robot's
development, covering gesture recognition accuracy, movement prediction performance,
system-level throughput, and a benchmarked comparison against two systems from the
literature review. Gesture recognition achieved 99.38\% test accuracy and 83--100\%
physical trial accuracy across six classes; movement prediction achieved an ADE of
25.8px and FDE of 32.4px and correctly recognized all five tested movement classes;
and the system sustained concurrent multi-subsystem operation within its
computational budget. These results, together with the limitations discussed in
Chapter 5, establish both what the system demonstrably achieves and what remains to
be independently validated -- most notably, end-to-end autonomous navigation and
larger-scale physical trials with structured per-trial logging.

\chapter{Conclusion and Recommendation}

\section{Summary of the Study}

This project set out to design and implement a robotic system capable of recognizing
human intention through hand gestures and body movement, and responding autonomously
without requiring cloud connectivity or GPU-class hardware. The system was built on a
Yahboom MicroROS-Pi5 platform, combining a MediaPipe HandLandmarker-based gesture
pipeline, a Multilayer Perceptron (MLP) gesture classifier trained on 19 engineered
geometric hand features, a Long Short-Term Memory (LSTM) network for short-horizon
movement prediction, and a spatial-zone-based multi-person tracking approach with
gesture-confirmation locking. Autonomous navigation infrastructure (SLAM Toolbox,
AMCL, Nav2) was integrated into the same pipeline to support movement between human
encounters.

Across the course of the project, six specific objectives were pursued: dataset
collection and preparation, gesture classifier development, movement-prediction model
development, integration of gesture and movement recognition with the robot's control
architecture, spatial-zone-based multi-person handling, and physical testing and
evaluation of the resulting system.

\section{Conclusion}

The gesture recognition subsystem achieved a test accuracy of 99.38\% on a balanced,
1,000-sample-per-class dataset of six gesture classes (Stop, Follow, Go, Left, Right,
Back), and physical trials across three participants performing each gesture ten times
returned per-gesture accuracy ranging from 83\% to 100\%, with five of the six gestures
staying within a narrow 7--12 percentage-point band. The movement prediction subsystem,
built on an LSTM trained on 152 motion sequences, achieved an Average Displacement
Error of 25.8px and a Final Displacement Error of 32.4px, and all five movement
classes (Approaching, Moving Away, Moving Left, Moving Right, Stationary) were
successfully and consistently recognized during functional testing.

Physical deployment surfaced and resolved several real engineering constraints not
apparent from simulation alone, including memory-pressure crashes on the original 2GB
Raspberry Pi 5 configuration (resolved by upgrading to 8GB), a \texttt{ROS\_DOMAIN\_ID}
mismatch silently preventing inter-node communication, and a brittle rule-based
gesture classifier that motivated the shift to a trained MLP. SLAM-based occupancy
grid mapping was successfully validated on physical hardware, producing a coherent map
of a real indoor environment. Autonomous path-planning through the resulting map via
Nav2 was architecturally integrated but had not been independently validated as of the
time of writing; this is addressed as a limitation below rather than claimed as an
achieved result.

Taken together, these results demonstrate that explicit, gesture-based human-robot
command interfaces can be combined with implicit, motion-based intention inference and
deployed on low-cost, edge-only hardware without cloud dependency — directly
addressing the gap identified in the literature review, where reviewed systems
achieved either strong intention prediction in isolation from navigation, or
navigation-integrated intention inference using only passive, inferred signals rather
than explicit human commands.

\section{Limitations}

Several limitations should be considered when interpreting the results of this study:

\begin{itemize}
  \item \textbf{Participant sample size.} Gesture recognition trials were conducted
  with three participants, each performing every gesture ten times (30 trials per
  gesture). This is sufficient to demonstrate functional viability but is too small a
  sample to draw strong statistical conclusions about generalization across a wider
  population, particularly for gestures such as Stop, which showed the widest
  accuracy spread (83--100\%).

  \item \textbf{Absence of raw per-trial records.} The gesture and movement testing
  results reported in Chapter 4 are drawn from summary-level data; individual per-trial
  outcomes were not logged in structured form at the time those trials were conducted.
  A dedicated data-logging node (\texttt{test\_logger\_node.py}) was developed during
  this project specifically to address this gap for future testing, but it was not yet
  in use for the results reported here.

  \item \textbf{Autonomous navigation not independently validated.} While SLAM-based
mapping has now been validated on physical hardware with drift-related issues
diagnosed and corrected (Section~\ref{sec:slam_reliability_fixes}), autonomous
path-planning through the resulting map using Nav2 had not been demonstrated
end-to-end at the time of writing. A configuration audit of the navigation stack
identified and corrected several issues preventing successful bringup, and the
specific remaining blocker has been identified and a fix prepared; however,
end-to-end autonomous navigation -- including obstacle avoidance, path-planning
reliability in a populated environment, and recovery behavior -- had not yet been
demonstrated successfully at the time of writing.

  \item \textbf{Map reusability and loop-closure tuning.} A systematic diagnosis
(Section~\ref{sec:slam_reliability_fixes}) identified and resolved two root
causes of map quality degradation -- a transform race condition causing high
scan-drop rates, and disabled yaw fusion causing rotational drift -- with both
fixes verified to produce a coherent, drift-free map. Two limitations remain
independent of these fixes, however: the occupancy grid map is saved only as a
flattened image/YAML export without a serialized SLAM state, meaning mapping
sessions cannot currently be resumed or incrementally improved and each new
session begins from a blank map; and loop-closure-specific parameters have not
been tuned beyond \texttt{slam\_toolbox}'s default configuration.

  \item \textbf{Facial recognition not deployed.} A facial recognition approach using
InsightFace (ArcFace embeddings) was evaluated offline, successfully distinguishing
between two enrolled operator identities with confidence scores of 0.54--0.57 against
an ``Unknown'' baseline of 0.09. A ROS2-integrated node
(\texttt{face\_id\_node}, designed against the deployed pipeline's existing
\texttt{/camera/image\_raw/compressed} and \texttt{/cognition/detection} topics) was
subsequently written toward productionising this capability, but it has not yet been
registered as a buildable package entry point, has not been compiled into the
workspace, and the InsightFace dependency is not installed in the robot's runtime
container. Facial recognition therefore remains unintegrated into the live system.

  \item \textbf{Person-following control stability.} The FOLLOW command's visual
servoing behaviour steers the robot's whole body (Section~\ref{sec:temporal_buffering},
Reactive Person Tracking) using a plain proportional controller on the horizontal
tracking error, which is susceptible to overshoot and oscillation around the target.
The camera itself is centred to a fixed pan/tilt position once at system startup and
does not track independently of the robot's body movement; a genuinely independent
pan-tilt tracking mechanism, if desired for improved facial-capture framing, has not
yet been implemented.

  \item \textbf{Hardware-specific validation.} All testing was conducted on a single
  physical unit under indoor lighting and environmental conditions typical of the test
  setting; performance under significantly different lighting, surface, or spatial
  conditions was not evaluated.
\end{itemize}

\section{Recommendations}

Based on the outcomes and limitations of this study, the following recommendations are
made for future extensions of this work:

\begin{itemize}
  \item \textbf{Expand physical trial scale and rigor.} Future testing should use the
  \texttt{test\_logger\_node.py} instrumentation developed in this project to capture
  structured, per-trial data across a larger and more diverse participant pool,
  enabling statistically robust accuracy, precision/recall, and confusion-matrix
  analysis rather than summary-level reporting.

  \item \textbf{Validate and tune autonomous navigation.} Nav2-based autonomous
  path-planning should be independently tested and validated end-to-end, including
  obstacle avoidance and recovery behavior, with results logged and reported
  separately from SLAM mapping validation.

  \item \textbf{Improve map persistence and quality.} Future work should implement
  serialized SLAM state saving to allow mapping sessions to be resumed and
  incrementally refined, and should tune \texttt{slam\_toolbox}'s loop-closure
  parameters -- ideally validated using a deliberate loop-closure driving pattern --
  to improve map consistency in larger or more complex environments.

  \item \textbf{Deploy and integrate facial recognition.} The offline facial
  recognition prototype should be integrated into the live perception pipeline on the
  robot, including resolving how it interacts with the existing gesture and
  spatial-zone detection logic for multi-operator disambiguation.

 \item \textbf{Upgrade person-following control to closed-loop PID.} Replacing the
current proportional-only visual-servoing controller with a tuned PID controller
(with Kalman filtering as a further refinement) would reduce tracking overshoot
during the FOLLOW behaviour. Independently, adding a dedicated pan-tilt tracking
mechanism -- decoupled from robot body movement -- would improve frame stability
specifically for facial recognition capture.

  \item \textbf{Broaden environmental and hardware validation.} Future testing should
  evaluate system performance across varied lighting conditions, floor surfaces, and
  physical layouts to better characterize the system's robustness outside the original
  test environment.
\end{itemize}
